\hypeexplica{../sections/hype-explica-z-vs-t/banner-hype-explica.png}
{Gustavo Neves Moraes}
{

\hspace*{1.5em}
Nos cursos introdutórios de estatística, o chamado \textit{z-test} e o \textit{t-test}, que costumam aparecer lado a lado,
com fórmulas praticamente idênticas e valores críticos distintos.
No entanto, a diferença entre eles não é apenas operacional,
pois está relacionada ao que se conhece sobre a população antes de a análise começar
e à incerteza que acompanha essa informação.

\subtitle{O teste de hipótese para a média}

\hspace*{1.5em}
Um teste de hipótese para a média responde a uma simples pergunta:
\highlight{A média observada na amostra, $\bar{x}$, é compatível com o valor de referência $\mu_0$, afirmado pela hipótese nula?}
Como duas amostras da mesma população raramente têm a mesma média,
alguma diferença entre $\bar{x}$ e $\mu_0$ é esperada mesmo quando $H_0$ é verdadeira.
O que o teste avalia é se essa diferença é pequena o bastante para ser atribuída ao acaso
ou grande demais para isso.

\hspace*{1.5em}
Para medir o tamanho da diferença, ela é expressa em unidades de erro padrão,
que indicam o quanto uma média amostral costuma variar de uma amostra para outra.
O resultado é a estatística de teste,
um número que informa a quantos erros padrão a média observada está de $\mu_0$.
Essa estatística é então comparada com um valor crítico, definido pelo nível de significância $\alpha$,
e, se ela o ultrapassa, rejeita-se $H_0$.

\hspace*{1.5em}
O valor crítico, por sua vez, depende da distribuição que a estatística segue quando $H_0$ é verdadeira.
É exatamente nesse ponto que o \textit{z-test} e o \textit{t-test} se separam.

\subtitle{Dois testes, uma diferença}

\hspace*{1.5em}
O erro padrão da média depende do desvio padrão populacional $\sigma$.
No \textit{z-test}, esse parâmetro é considerado conhecido,
e a estatística de teste segue a distribuição normal padrão:

\[
z = \frac{\bar{x} - \mu_0}{\sigma / \sqrt{n}}
\]

\hspace*{1.5em}
No \textit{t-test}, $\sigma$ é desconhecido e substituído pelo desvio padrão amostral $s$.
A fórmula tem a mesma aparência, mas a estatística passa a seguir a distribuição t de Student,
com $n - 1$ graus de liberdade:

\[
t = \frac{\bar{x} - \mu_0}{s / \sqrt{n}}
\]

\begin{hypeitemize}
    \item \textbf{\textit{z-test}:} $\sigma$ conhecido, distribuição normal padrão.
    \item \textbf{\textit{t-test}:} $\sigma$ estimado por $s$, distribuição t de Student.
\end{hypeitemize}

\subtitle{Por que as caudas são mais pesadas?}

\hspace*{1.5em}
No \textit{z-test}, o único elemento que varia de uma amostra para outra é a média $\bar{x}$,
pois o desvio padrão $\sigma$ é um valor fixo e conhecido.
No \textit{t-test}, o desvio padrão amostral $s$ também é calculado a partir dos dados
e, portanto, muda a cada amostra.
Isso significa que o denominador da estatística passa a ser uma fonte extra de variação.

\hspace*{1.5em}
Essa variação adicional importa porque $s$ às vezes subestima $\sigma$ por puro acaso.
Quando isso acontece, o erro padrão calculado fica menor do que deveria,
e a estatística resultante se afasta de zero mais do que o esperado,
mesmo com $H_0$ verdadeira.
Em amostras pequenas, $s$ oscila bastante,
então esses valores extremos aparecem com mais frequência do que a distribuição normal prevê.

\hspace*{1.5em}
A distribuição t de Student incorpora esse efeito por meio de caudas mais pesadas,
ou seja, atribui mais probabilidade aos valores extremos.
Como consequência, os valores críticos são maiores,
e o critério para rejeitar $H_0$ se torna mais exigente,
o que compensa a incerteza de ter estimado $\sigma$ a partir de poucos dados.

% \articleimage{1.0}{../sections/hype-explica-z-vs-t/distribution.png}{AA}

\subtitle{Graus de liberdade e convergência}

\hspace*{1.5em}
O formato da distribuição t é determinado pelos graus de liberdade,
que no teste para uma única média valem $n - 1$.
Com poucos graus de liberdade, as caudas se afastam bastante das da normal,
e, à medida que $n$ aumenta, $s$ passa a estimar $\sigma$ com mais precisão,
de modo que a distribuição t converge para a normal padrão.
A figura e a tabela a seguir ilustram esse comportamento.

\articleimage{1.0}{../sections/hype-explica-z-vs-t/curvas-z-t.png}{Distribuição normal padrão comparada a distribuições t com 2, 5 e 30 graus de liberdade.}

\resizebox{\linewidth}{!}{%
\begin{tabular}{|cccc|}
\hline
\textbf{\small Amostra ($n$)} & \textbf{Graus de liberdade} & \textbf{Valor crítico t} & \textbf{Valor crítico z} \\
\hline
5   & 4   & 2,776 & 1,960 \\
10  & 9   & 2,262 & 1,960 \\
30  & 29  & 2,045 & 1,960 \\
100 & 99  & 1,984 & 1,960 \\
\hline
\end{tabular}%
}

\vspace{0.4em}
{\small Valores críticos bilaterais para $\alpha = 0{,}05$.}

\hspace*{1.5em}
Para $n = 5$, o valor crítico t é cerca de 40\% maior que o z,
enquanto para $n = 100$ a diferença é inferior a 2\%.
A distância entre os dois testes, portanto, é relevante principalmente em amostras pequenas.

\subtitle{Um exemplo numérico}

\hspace*{1.5em}
Considere uma amostra de $n = 5$ observações, com média $\bar{x} = 108$ e desvio padrão $s = 8$,
usada para testar $H_0: \mu = 100$ com $\alpha = 0{,}05$.
O erro padrão é $8/\sqrt{5} \approx 3{,}58$,
e a estatística de teste vale $(108 - 100)/3{,}58 \approx 2{,}24$.

\hspace*{1.5em}
Se $s$ for tratado como se fosse $\sigma$ e comparado ao valor crítico z de 1,96,
a conclusão é rejeitar $H_0$, com p-valor em torno de 0,025.
Ao utilizar o valor crítico t para 4 graus de liberdade, que é 2,776,
não há evidência suficiente para rejeitar $H_0$, e o p-valor passa a ser aproximadamente 0,09.
Os mesmos dados conduzem a conclusões opostas,
e apenas a segunda leva em conta que o desvio padrão foi estimado a partir de cinco observações.

\vspace{12mm}
\subtitle{Quando cada teste se aplica}

\hspace*{1.5em}
O \textit{z-test} é adequado quando $\sigma$ é de fato conhecido,
situação rara, que costuma ocorrer em processos industriais com longo histórico de calibração.
Ele também é o procedimento padrão para proporções,
caso em que a variância é determinada pela própria proporção testada
e não precisa ser estimada separadamente.
Para médias com $\sigma$ desconhecido, que é o cenário usual em ciência de dados,
o \textit{t-test} é o procedimento apropriado.

\subtitle{Uma ressalva importante}

\hspace*{1.5em}
A regra segundo a qual $n \geq 30$ autoriza o uso do \textit{z-test} é uma convenção histórica,
associada à época em que as tabelas eram impressas e consultadas manualmente.
Com $n$ grande, a diferença entre os dois testes se torna desprezível,
e, como os softwares estatísticos calculam o \textit{t-test} sem custo adicional,
a prática atual é utilizá-lo sempre que $\sigma$ for desconhecido.
Vale lembrar ainda que, em amostras pequenas, o \textit{t-test} pressupõe uma população aproximadamente normal,
condição que deve ser verificada antes de se confiar no resultado.

\begin{infobox}
\textbf{Curiosidade:}\\
A distribuição t foi proposta por William Sealy Gosset em 1908,
quando trabalhava na cervejaria Guinness, em Dublin,
e a empresa o obrigou a publicar sob o pseudônimo \textit{Student}.
Ele precisava tirar conclusões a partir de lotes pequenos de cevada e de cerveja,
um contexto em que as aproximações normais não funcionavam bem.\\
\textbf{\href{https://www.instagram.com/p/Dexi6WNFQJu/}{link do vídeo sobre a curiosidade aqui.}}
\end{infobox}

\vspace{0.1em}
{\color{HypePurple2} Quer se aprofundar mais?}

\href{https://www.open.edu/openlearn/science-maths-technology/data-analysis-hypothesis-testing/content-section-6.1}{- Data analysis: hypothesis testing} \\
\href{https://www.jstor.org/stable/2331554?origin=crossref}{- Student (1908), The Probable Error of a Mean, Biometrika, vol. 6.}

}